% ===== 原文 pp.24–26 =====
% （§4.2 的后半 + §4.3、§4.3.1、§4.3.2、§4.3.3 的开头；本块 tag 4.8–4.21）

% [原文 p.24]
相应地，对 $\cN=2$ 拉格朗日量场论，$Z(q,g)$ 的计算既可以通过在平坦空间上数算符来完成，也可以通过做 $S^1\times S^3$ 上的路径积分来完成。文献中大多用算符方法；对 $\cN=2$ 情形的透彻讨论见 [9]。对 $\cN=4$ 的路径积分方法，见 [25]。

% [Tutor Note: p.24, 原文此处有两处排印笔误（已用 700 dpi 裁图复核确认原文如此）：
%  "In most of th literature" 应为 "In most of the literature"；"for a through discussion" 应为
%  "for a thorough discussion"。译文按正确含义处理，公式与结论不受影响。]

无论用哪种方式，我们都得到如下结果。首先把绕 $S^1$ 的和乐 $g\in G$ 用复数写成
\begin{equation}
g=(z_1,\ldots,z_r)\in\U(1)^r\subset G
\tag{4.8}
\end{equation}
的形式，其中 $r$ 是 $G$ 的秩。味和乐可以任意取。规范和乐则参数化我们对其积分的 BPS 组态。

% [Tutor Note: p.24, (4.8)–(4.10) 的分工务必分清：味和乐 $g$ 是「可任意取」的背景参数
%  （它把质量参数/味化学势代进来），而规范和乐才是「要被积分」的变量——(4.10) 的围道积分
%  $\oint\prod_i\dd z_i/(2\pi\sqrt{-1}z_i)$ 就是对这个规范和乐取的，且按原文「在原点取留数」。
%  原文 p.24 那句 "The gauge holonomy parameterizes the BPS configurations over which we integrate"
%  正是这个意思；把规范、味两种和乐弄反，是本节最典型的错误。]

对一个由处于对称性 $G$ 的表示 $R$ 中的 $\cN=1$ 手征多重态（chiral multiplet）所构成的 $\cN=2$ 超多重态，配分函数为
\begin{equation}
Z_R(q,g)=\prod_{n\ge0}\prod_w\frac{1}{1-q^{n+1/2}z^w}
\tag{4.9}
\end{equation}
其中 $w=(w_1,\ldots,w_r)$ 跑遍 $R$ 的权，而 $z^w:=\prod z_i^{w_i}$。

对一个规范群为 $G$ 的规范理论，配分函数为
\begin{equation}
Z(q)=\frac{1}{|W_G|}\oint\prod_{i=1}^r\frac{\dd z_i}{2\pi\sqrt{-1}z_i}\prod_\alpha(1-z^\alpha)K(z)^{-2}\times(\text{物质贡献})
\tag{4.10}
\end{equation}
其中
\begin{equation}
K(z)^{-1}=\prod_{n\ge0}\left[(1-q^{n+1})^r\prod_\alpha(1-q^{n+1}z^\alpha)\right].
\tag{4.11}
\end{equation}
这里对 $\alpha$ 的乘积跑遍 $G$ 的根 $\alpha$，$|W_G|$ 是 $G$ 的 Weyl 群的阶，而积分取原点处的留数。

% [Tutor Note: p.24, (4.11) 中 $(1-q^{n+1})^r$ 的 $r$ 是规范群 $G$ 的秩，它与后面 $\prod_\alpha$
%  的根因子合起来恰好给出 $\dim G$ 个「单粒子字母」：以 $\SU(N)$ 为例，$r=N-1$、正负根共 $N(N-1)$ 个，
%  $(N-1)+N(N-1)=N^2-1=\dim\SU(N)$。所以 $(1-q^{n+1})^r$ 不是笔误漏掉的多余指数，不要擅自改掉。]

% [Tutor Note: p.24, (4.9) 中 $q$ 的指数是 $n+\tfrac12$、求和从 $n=0$ 起，故 $q^{1/2}$ 这样的分数幂
%  是正常的（它来自 $\cN=1$ 手征多重态的 R 荷/共形维数组合），不是排版错误；把它误读成 $n+1$ 会
%  使整个指标的 $q$-展开错位。另注意 (4.9) 只是「一个」手征多重态的贡献，$\cN=2$ 超多重态的完整
%  贡献要把两半相乘——这正是 (4.12) 分母里出现 $\prod_\pm$ 的原因。]

作为例子，规范群为 $\SU(N)$ 的 $\cN=4$ 超杨–米尔斯理论，若看作具有 $\SU(2)$ 味对称性的 $\cN=2$ 理论，其配分函数给出如下。$\SU(2)$ 的元素 $g$ 写成 $(a,a^{-1})\in\SU(2)$。$\SU(N)$ 的元素写成 $(z_1,\ldots,z_N)\in\SU(N)$，并约定 $\prod z_i=1$。于是我们有
\begin{equation}
\begin{aligned}
Z(q,a)&=\frac{1}{N!}\oint\prod_{i=1}^{N-1}\frac{\dd z_i}{2\pi\sqrt{-1}z_i}\prod_{i,j}(1-z_i/z_j)^2\\
&\quad\times\frac{\prod_{n\ge0}(1-q^{n+1})^{2r}\prod_{i,j}(1-q^{n+1}z_i/z_j)^2}{\prod_\pm\prod_{n\ge0}(1-q^{n+1/2}a^{\pm1})^r\prod_{i,j}(1-q^{n+1/2}z_i/z_j a^{\pm1})}.
\end{aligned}
\tag{4.12}
\end{equation}

% [Tutor Note: p.24, (4.12) 是 (4.10)(4.11) 在 $G=\SU(N)$、物质为伴随超多重态时的显式化：
%  分子中的 $\prod_{i,j}(1-q^{n+1}z_i/z_j)^2$ 来自矢量多重态（把伴随表示的特征标写成 $z_i/z_j$），
%  分母中的 $\prod_\pm$ 与 $q^{n+1/2}$ 来自伴随超多重态（味 $\SU(2)$ 二重态给出 $a$ 与 $a^{-1}$ 两份）。
%  核对时注意 $\prod_{i,j}$ 已自动包含 $i=j$ 的项（此时 $z_i/z_j=1$）。]

% [Tutor Note: p.24, 注意 (4.12) 第一行的 $\prod_{i,j}(1-z_i/z_j)^2$ 若真按 $i,j=1,\dots,N$ 读，
%  $i=j$ 时因子恰为 $(1-1)^2=0$，整个被积函数会变成零。与 (4.10) 的 $\prod_\alpha(1-z^\alpha)$ 对照可知
%  这里应是 $\prod_{i<j}$（或 $\prod_{i\ne j}$），其平方正是伴随表示 Weyl 因子的平方。
%  已用 700 dpi 裁图复核：原文确实只印了下标 $i,j$、没有写 $i<j$，属原文的记号疏漏；正文照原样排印。
%  注意分子里 $\prod_{i,j}(1-q^{n+1}z_i/z_j)^2$ 的 $i=j$ 因子为 $(1-q^{n+1})^2\ne0$，那里没有这个问题。
%  同样的疏漏见 p.28 的 (4.27)。]

这个表达式看起来很复杂，但它相当显式，而且它的 $q$-展开可以很容易地算出。

\subsection{$S^4$}

作为第二个超对称背景，我们考虑 $S^4$。圆 $S^4$ 上的局域化最早在 [6] 中完成，后来在 [26] 中被推广到压扁的 $S^4$。我们把这两种背景都简称为 $S^4$。关于一般双等距背景上局域化的综述见 [27]。

% [原文 p.25]
\subsubsection{一般结构}

为描述局域化，我们首先把 $S^4$ 写成五维空间中的超曲面，该空间由实参数 $x$ 与两个复数 $z_{1,2}$ 参数化：
\begin{equation}
x^2+|z_1|^2+|z_2|^2=1.
\tag{4.13}
\end{equation}
记 $z_i=r_i\ee^{\sqrt{-1}\phi_i}$，我们就给 $S^4$ 赋以在 $\phi_1$ 与 $\phi_1$ 的任意平移下不变的背景超引力场。

% [Tutor Note: p.25, 原文 "invariant under the arbitrary shift of φ_1 and φ_1" 中两个下标都印成了 1，
%  已用 700 dpi 裁图复核确认原文如此；按上下文应为 "$\phi_1$ 与 $\phi_2$"（即两个复坐标 $z_{1,2}$ 的辐角）。
%  正文照原样排印，此处按正确含义理解。]

度规与其它超引力场可以适当地选取，使得存在一个超等距变换 $\epsilon$，其 $\delta_\epsilon^2$ 生成一个转动
\begin{equation}
\phi_i\to\phi_i+\epsilon_i.
\tag{4.14}
\end{equation}
注意 $x=1,z_i=0$ 与 $x=-1,z_i=0$ 是这个转动下的两个不动点。我们称它们为北极与南极。

% [Tutor Note: p.25, (4.13)–(4.14) 的几何：$S^4$ 是 $\C^2\times\R$ 中由
%  $x^2+|z_1|^2+|z_2|^2=1$ 给出的超曲面，$\phi_i\to\phi_i+\epsilon_i$ 是 $\U(1)^2\subset\SO(5)$ 的等距变换，
%  北极/南极即 $x=\pm1$ 的两个不动点。$\delta_\epsilon^2$ 生成这个等距变换，正是它保证局域化轨迹
%  只在南北极附近有支撑——这是后面 (4.15) 能把路径积分化为对 $a$ 的有限维积分、且结果呈
%  $Z(a)\overline{Z(a)}$ 两个半球贡献之积的原因。]

在局域化下对配分函数有贡献的场组态如下。对超多重态，真空期望值（vev）应全为零。对矢量多重态，伴随标量 $\Phi$ 使得 $\Phi$ 可以有非零的、与时空无关的真空期望值，但 $\bar\Phi=0$。这个真空期望值习惯上记作 $a$。我们取 $a$ 为纯虚数的约定。当矢量多重态是非动力学的时候，这个真空期望值 $a$ 给出超多重态的质量项，并常记作 $m$。此外，还可以有支撑在北极的点状瞬子与支撑在南极的点状反瞬子。

相应地，$S^4$ 上规范群为 $G$ 的理论的配分函数具有如下形式
\begin{equation}
Z=\frac{1}{|W_G|}\int\dd^r a\,Z(a)\overline{Z(a)}
\tag{4.15}
\end{equation}
其中 $|W_G|$ 是 Weyl 群的阶，积分跑遍 $\Phi$ 实部的真空期望值空间（已规范固定到位于 Cartan 子代数上）。$Z(a)$ 是给定真空期望值 $a$ 时来自北半球的贡献，其复共轭 $\overline{Z(a)}$ 给出南半球的贡献。

贡献 $Z(a)$ 常被称为 Nekrasov 配分函数，它由以下成分组成：
\begin{equation}
Z(a)=Z_{\text{classical}}(a)Z_{\text{one-loop}}(a)Z_{\text{instanton}}(a).
\tag{4.16}
\end{equation}
下面我们给出这三个因子的显式形式。对于与伴随标量的真空期望值、质量参数和规范耦合常数无关的整体因子，我们将不作深究。

% [Tutor Note: p.25, (4.15) 的因子 $1/|W_G|$ 是 Weyl 群的阶；对 $G=\SU(N)$ 有 $|W_G|=N!$，
%  所以后面 (4.25) 里出现的 $1/N!$ 与 (4.15) 的 $1/|W_G|$ 是同一个归一化因子，并非两套不同的约定。
%  另外被积函数是 $Z(a)\overline{Z(a)}=|Z(a)|^2$：北极贡献与南极贡献互为复共轭，这正是 $S^4$ 上
%  局域化结果呈「模平方」形状的原因，也是 (4.25) 与 (4.27) 里那些 $|\cdot|^2$ 的来源。]

\subsubsection{经典因子与单圈因子}

第一个因子是经典作用量的指数，它由下式在各个规范多重态上的乘积给出：
\begin{equation}
\ee^{-\frac{1}{\epsilon_1\epsilon_2}2\pi\sqrt{-1}\tau\langle a,a\rangle}
\tag{4.17}
\end{equation}
% [原文 p.26]
其中 $\tau$ 是复化规范耦合。

% [Tutor Note: p.25–26, (4.17) 的结构：指数上的 $2\pi\sqrt{-1}\tau$ 是复化耦合，$\langle a,a\rangle$
%  是 Cartan 上的二次型（在 $\SU(N)$ 中即 $\sum_i a_i^2$，这正是 (4.25)(4.27) 中 $\sum a_i^2$ 的来源），
%  而 $1/(\epsilon_1\epsilon_2)$ 是 $\Omega$ 背景的体积因子。注意 (4.17) 只是「一个」规范多重态的因子，
%  原文紧接着说要对「各个规范多重态」取乘积；且 (4.16) 把经典因子定义成整个经典作用量的指数，
%  所以 (4.17) 是逐因子的写法，不要漏掉对规范群各简单因子的乘积。]

第二个因子是来自矢量多重态与超多重态的单圈贡献。来自规范多重态，我们有
\begin{equation}
Z^G_{\text{one-loop}}=\prod_{\alpha>0}\frac{1}{\Gamma_2(\alpha\cdot a+\epsilon_1+\epsilon_2|\epsilon_1,\epsilon_2)\Gamma_2(\alpha\cdot a|\epsilon_1,\epsilon_2)}
\tag{4.18}
\end{equation}
其中乘积跑遍正根 $\alpha$。来自处于表示 $R$ 的超多重态，我们有
\begin{equation}
Z^R_{\text{one-loop}}=\prod_w\Gamma_2\Big(w\cdot a+\frac{\epsilon_1+\epsilon_2}{2}\Big|\epsilon_1,\epsilon_2\Big)
\tag{4.19}
\end{equation}
其中乘积跑遍 $R$ 的权。对 $R$ 中的半超多重态，乘积跑遍 $R$ 的权中任意一半；最终对瞬子积分 (4.15) 的贡献并不依赖于这种对半的划分。

% [Tutor Note: p.26, $\Gamma_2(x|\epsilon_1,\epsilon_2)$ 是 Barnes 双 Gamma 函数，下标 2 指「双」，
%  不是平方，不要读成 $\Gamma^2$。本综述没有给出它的定义，按 Nekrasov 文献的常用约定
%  $\Gamma_2(x|\epsilon_1,\epsilon_2)=\prod_{m,n\ge0}(x+m\epsilon_1+n\epsilon_2)^{-1}$ 看结构即可：
%  (4.18) 矢量多重态贡献的是它的倒数（出现在分母），(4.19) 超多重态贡献的是它本身（出现在分子）。
%  易错点：为什么 (4.18) 里 $\alpha\cdot a+\epsilon_1+\epsilon_2$ 后面没有 $1/2$、而 (4.19) 里却有 $(\epsilon_1+\epsilon_2)/2$？
%  这是矢量与超多重态 R 荷分配不同的结果，不是笔误——用 (4.26) 的化简做反查即可确认。]

超多重态的质量项可以通过给味背景规范多重态中的标量一个真空期望值来纳入。有效地说，这只不过把 $w\cdot a\to w\cdot a+m$。还要注意，在我们的约定中 $a$ 是纯虚数而 $\epsilon_{1,2}$ 是实数。因此，为得到来自南半球的贡献 $\overline{Z(a)}$，我们只需在双 Gamma 函数的自变量中把 $a$ 换成 $-a$。

\subsubsection{瞬子贡献}

瞬子贡献的表述要复杂得多，而且由于各种技术问题，其显式形式只对 SU 规范群、且只有完整超多重态时才已知。即使对 SU 规范群，计算也牵涉到某种会引入虚假贡献的正规化；人们常说这些虚假贡献来自把 SU 群换成 U 群，因而有「U(1) 因子」之名。这里我们只引用已知结果，其中包含虚假因子或 U(1) 因子。更详细的讨论见 [28]。

% [Tutor Note: p.26, 这里的限制是「显式公式」的限制，而不是理论本身的限制：原文说的是
%  "the explicit form is only known for SU gauge groups and with full hypermultiplets"。
%  p.27 给出的 (4.24) 是双基表示超多重态的情形，而伴随表示与基础表示都只是它的特化：
%  令 $\vec a=\vec b,\vec Y=\vec W$ 得到伴随表示；令 $M=N_f$、$W_i=0$ 并把 $\SU(M)$ 当作背景规范场
%  得到 $N_f$ 个基础表示。用这个对应关系可以自查 (4.24) 的每一处下标。]

一个 $\U(N)$ 规范多重态的点状瞬子组态由一个 $N$ 元组 Young 图 $\vec Y=(Y_1,\ldots,Y_N)$ 标记。一个 Young 图 $Y$ 只是一个不增的自然数序列 $Y=(\lambda_1\ge\lambda_2\ge\cdots)$。这里我们把它可视化为：把 $\lambda_i$ 看作从左数第 $i$ 列的高度，并当 $i$ 大于图的宽度时令 $\lambda_i=0$。这不是 Young 图文献中的标准约定，但却是瞬子计算中的标准约定。我们用 $|Y|$ 表示盒子数，即 $|Y|=\sum\lambda_i$，然后定义 $|\vec Y|=\sum|Y_i|$。

% [Tutor Note: p.26, 原文这句写作 "a non-decreasing sequence of natural numbers Y=(λ1≥λ2≥…)":
%  词（non-decreasing，非降）与式子（$\lambda_1\ge\lambda_2\ge\cdots$，非增）矛盾；按 Young 图的
%  约定与紧随其后的画法，正确者应是「非增」。已用 700 dpi 裁图复核原文确实印作 "non-decreasing"。
%  译文取「不增」以免读者误导，此处仅作记录。]

% [Tutor Note: p.26, Young 图的约定是本节最易与教科书相反的地方：这里 $\lambda_i$ 是「第 $i$ 列的高度」
%  （教科书通常把 $\lambda_i$ 当第 $i$ 行的长度），所以 $Y=(5,4,3,2,2,1)$ 画出来是图 11 的阶梯形；
%  $\lambda'_j$ 则是第 $j$ 行的长度。把 $s=(i,j)$ 理解为「第 $i$ 列、第 $j$ 行」（行自下而上数）后：
%  $A_Y(s)=\lambda_i-j$ 数的是同一列中 $s$ 上方的盒子数，$L_Y(s)=\lambda'_j-i$ 数的是同一行中
%  $s$ 右侧的盒子数。以图 11 的 $s=(2,2)$ 核对：$\lambda_2=4\Rightarrow A_Y=4-2=2$（$s$ 上方 2 个空心圆），
%  $\lambda'_2=5\Rightarrow L_Y=5-2=3$（$s$ 右侧 3 个实心点）。关键：$s$ 允许在 $Y$ 之外，
%  此时 $A_Y(s)$ 或 $L_Y(s)$ 可取负值（例如紧贴图右侧的盒子有 $A_Y=-1$），这是 (4.23) 的乘积
%  能写对的前提，不是笔误。]

我们用 $Y^T$ 表示它的转置，$Y^T=(\lambda'_1\ge\lambda'_2\ge\cdots)$。对位于坐标 $(i,j)$ 的盒子 $s$，其臂长 $A_Y(s)$ 与腿长 $L_Y(s)$ 定义为
\begin{equation}
A_Y(s)=\lambda_i-j,\qquad L_Y(s)=\lambda'_j-i.
\tag{4.20}
\end{equation}
注意 $s$ 可以位于 Young 图 $Y$ 之外。于是我们令
\begin{equation}
E(a,Y_1,Y_2,s)=a-\epsilon_1L_{Y_2}(s)+\epsilon_2(A_{Y_1}(s)+1).
\tag{4.21}
\end{equation}

% [Tutor Note: p.26, (4.21) 里两个下标的交叉最易抄错：$\epsilon_1$ 配的是 $L_{Y_2}$（腿长属于第二个图），
%  $\epsilon_2$ 配的是 $(A_{Y_1}+1)$（臂长属于第一个图）。检查办法：看 (4.23)(4.24) 中
%  $E(a_i-a_j,Y_i,Y_j,s)$ 的写法——$Y_i$ 配臂、$Y_j$ 配腿，与 (4.21) 的「$Y_1$ 配 $A$、$Y_2$ 配 $L$」一致。
%  把 $Y_1\leftrightarrow Y_2$ 或 $\epsilon_1\leftrightarrow\epsilon_2$ 弄混，(4.25) 的化简就无法与 (4.26) 对上。]
